% !TEX root = ../main.tex
%--- Slide 11: conclusioni (circa 1 min) --------------------------------------------------
\begin{frame}{Conclusioni}
  \begin{itemize}\setlength{\itemsep}{3pt}
    \item Nei dati reali DM e scattering sono \textbf{degeneri} (due picchi, forma a L);
          nelle simulazioni senza processi non modellati no
    \item \textbf{GWB} non modellato: solo deboli code a L
    \item \textbf{Band noise} non modellato: sposta il DM e ne riproduce il picco ad alta ampiezza
          dei dati reali, ma non la forma a L
    \item Un band noise può indurre un bias nel DM senza essere rivelato:
          \textbf{non può essere esclusa} la sua presenza nei dati reali
    \item Nessuna simulazione riproduce tutta la struttura:
          l'origine della degenerazione \textbf{resta aperta}
  \end{itemize}
  \vspace{2mm}
  \centering
  \fbox{\parbox{0.85\textwidth}{\centering
    Un processo non modellato può distorcere tutte le altre componenti:\\
    un modello di rumore accurato è essenziale per la ricerca del GWB}}
  \vspace{1mm}
  \begin{flushleft}
  {\small\color{grigio} Sviluppi: altri processi non modellati, altre realizzazioni del rumore,
   altre pulsar di DR2low}
  \end{flushleft}
\end{frame}
